% sections/morse.tex -- the Morse order: Lemma 2 and Corollaries 2.1--2.2. Uses macros.tex. Executable checks of every statement:
% research/mmhm/reference/check_props.py.
\section{The Morse order}\label{sec:morse}

Discrete Morse theory~\cite{forman2002user} enters only through the choice of
keys at a rebuild. A rebuild computes an acyclic matching by collapses, in the
spirit of~\cite{benedetti2014random}, and orders cells by reverse removal time.

\begin{definition}[Collapse matching and Morse order]\label{def:collapse}
Run the following on $P:=K$, with a counter $t:=0$ and a queue $Q$ of the cells
$x$ that have exactly one immediate coface in $P$ (\emph{free faces}).
\begin{enumerate}
\item While $Q\ne\emptyset$: pop $x$; skip it unless $x\in P$ and $x$ has
  exactly one immediate coface $y$ in $P$. Then \emph{pair} $(x,y)$, set
  $\rmt(y):=t$, $\rmt(x):=t+1$, $t\gets t+2$, remove $y$ and $x$ from $P$, and
  push every facet of $y$ or $x$ that has become free.
\item If $P\ne\emptyset$, choose a cell $c$ of the highest dimension present in
  $P$ (smallest vertex tuple), declare it \emph{critical}, set $\rmt(c):=t$,
  $t\gets t+1$, remove it, push the facets that have become free, and go to~1.
\end{enumerate}
The \emph{Morse order} is $\ord(x):=-\rmt(x)$ (relabelled to $0,\dots,|K|-1$).
Write $M$ for the set of pairs, $\Crit$ for the critical cells, $\Lo$ for the
cells $x$ with $(x,y)\in M$ and $\Up$ for the cells $y$ with $(x,y)\in M$; for
$x\in\Lo$ let $\mu(x)$ be its partner. For a cell $x$, its \emph{youngest facet}
is the facet with the largest key and its \emph{oldest coface} the coface with
the smallest key.
\end{definition}

The removed cells are always maximal in $P$, and $P$ remains a subcomplex. For
the paired cell $y$ this uses the simplicial diamond property: if $y$ had a
coface $w\in P$, then $w$ would have exactly two facets containing $x$, namely
$y$ and some $y'\ne y$, and $y'\in P$ because $P$ is a subcomplex, so $x$ would
not be free. The queue discipline is not specified; all results below hold for
every discipline. \Cref{fig:morse} runs the procedure on a small complex.

A \emph{gradient path} (V-path) is a sequence $\tau_0,\sigma_0,\tau_1,\sigma_1,
\dots,\tau_m$ with $(\tau_i,\sigma_i)\in M$ and $\tau_{i+1}$ a facet of
$\sigma_i$ different from $\tau_i$.

\begin{lemma}[Gradient pairs are apparent pairs]\label{lem:apparent}
For every execution of the collapse procedure, with $\ord=-\rmt$:
\begin{enumerate}
\item[(a)] $\ord$ is a filtration order: $\ord(f)<\ord(s)$ whenever $f$ is a facet
  of $s$.
\item[(b)] $\rmt$ strictly increases along the lower cells of every gradient
  path ($\rmt(\tau_{i+1})>\rmt(\tau_i)$); in particular $M$ has no closed
  gradient path, so it is an acyclic matching.
\item[(c)] For every $(\tau,\sigma)\in M$: $\tau$ is the youngest facet of
  $\sigma$, i.e.\ $\tau=\low(\bd\sigma)$, and $\sigma$ is the oldest coface of
  $\tau$.
\item[(d)] If $z\ne0$ is a cycle, then $\low(z)\notin\Up$.
\item[(e)] In the static reduction under $\ord$, with or without clearing, for
  every $(\tau,\sigma)\in M$ the column $\sigma$ is not cleared, receives no
  addition, and obtains pivot $\tau$; thus $R_\sigma=\bd\sigma$ and
  $(\tau,\sigma)$ is an \emph{apparent pair}~\cite{bauer2021ripser}.
\end{enumerate}
\end{lemma}

\begin{proof}
(a) Only maximal cells of $P$ are removed (above), so every cell is removed
before each of its facets.

(b) When $\sigma_i$ is removed at time $\rmt(\sigma_i)$, its facet $\tau_{i+1}$ is
still in $P$. Times $\rmt(\sigma_i)$ and $\rmt(\tau_i)=\rmt(\sigma_i)+1$ are
used by $\sigma_i$ and $\tau_i$, and $\tau_{i+1}\ne\tau_i$, so
$\rmt(\tau_{i+1})>\rmt(\tau_i)$. A closed gradient path would give
$\rmt(\tau_0)<\rmt(\tau_0)$.

(c) The pair is removed at consecutive times, so $\ord(\tau)=\ord(\sigma)-1$:
the two cells are adjacent in the order. By (a) every facet of $\sigma$ precedes
$\sigma$, hence has key at most $\ord(\tau)$, and every coface of $\tau$ follows
$\tau$, hence has key at least $\ord(\sigma)$. (This is the elementary fact that
a facet and a coface that are adjacent in a filtration order form an apparent
pair.)

(d) Let $u=\low(z)\in\Up$ with $u=\mu(w)$. The coefficient of $w$ in $\bd z=0$
vanishes and $w\in\bd u$, so some $u'\in z$ with $u'\ne u$ has $w$ as a facet.
By (c), $u$ is the oldest coface of $w$, so $\ord(u')>\ord(u)$, contradicting
$u=\low(z)$.

(e) If $\sigma$ were cleared, it would be the low of a reduced column $R$ of the
next dimension; $R=\bd V'$ (Proposition~\ref{cor:rebuild}) is a non-zero cycle with
low $\sigma\in\Up$, contradicting (d). Columns are processed in increasing key
order. When $\sigma$ is processed, every earlier column $x$ has $R_x=0$ (if
cleared) or $R_x=\bd V_x$ with all cells of $V_x$ of key at most
$\ord(x)<\ord(\sigma)$ (Proposition~\ref{cor:rebuild}); by (c) no such cell has
$\tau$ as a facet, so $\tau\notin R_x$ and $\pivot(\tau)$ is undefined. Since
$\low(\bd\sigma)=\tau$ by (c), the reduction loop performs no addition and
assigns $\pivot(\tau)=\sigma$.
\end{proof}

\begin{corollary}[Work happens only on critical columns]\label{cor:critical-work}
In the static reduction with clearing under the Morse order, the column of every
$x\in\Lo$ with $\dim x\ge1$ is cleared, the column of every $y\in\Up$ is
apparent, and therefore every column addition is performed while reducing the
column of a critical cell.
\end{corollary}

\begin{proof}
For $(x,y)\in M$ with $\dim y=k+1\ge2$, Lemma~\ref{lem:apparent}(e) applied to
$\Bm{k+1}$, which is processed before $\Bm{k}$, gives $R_y=\bd y$ with
$\low(R_y)=x$, so $x$ is cleared. Upper cells are apparent by
Lemma~\ref{lem:apparent}(e). Every remaining column belongs to a critical cell.
\end{proof}

\begin{figure}[t]
\centering
% figures/f3_morse.tex -- the Morse order on a small complex: collapse, gradient
% pairs, critical cells, apparent pairs in the boundary matrices, and the
% reduction of the critical column. Rows and columns of every matrix are in key
% order (rows top to bottom), so a low is the lowest non-zero entry.
% Every entry is computed and checked by reference/figure_examples.py.
\newcommand{\pv}[1]{{\setlength{\fboxsep}{0.6pt}\setlength{\fboxrule}{0.6pt}\fbox{$#1$}}}
\newcommand{\cc}{\cellcolor{black!14}}
\newcommand{\gz}[1]{\textcolor{black!40}{#1}}
{\footnotesize\renewcommand{\arraystretch}{1.06}%
%
% ====================== (a) complex and matching ==========================
\begin{minipage}[t]{0.29\textwidth}
\textbf{(a)} Complex, matching, critical cells\par\medskip
\centering
\begin{tikzpicture}[>={Stealth[length=2.1mm,width=1.8mm]},
    mvtx/.style={circle, draw=black!85, fill=white, inner sep=0pt, minimum size=3.8mm,
      font=\scriptsize},
    cvtx/.style={mvtx, fill=black!22, line width=1.1pt},
    medg/.style={black!60, semithick},
    cedg/.style={black, line width=1.7pt},
    pair/.style={->, black, thick},
    lab/.style={font=\scriptsize, text=black!70, inner sep=1.5pt}]
  \coordinate (p0) at (0,0);     \coordinate (p1) at (2.1,0);
  \coordinate (p2) at (2.1,2.1); \coordinate (p3) at (0,2.1);
  \fill[black!10] (p0) -- (p1) -- (p2) -- cycle;
  \draw[medg] (p0) -- (p1) -- (p2) -- (p3) -- cycle;
  \draw[cedg] (p0) -- (p2);
  \draw[pair] ($(p0)+(0,0.22)$) -- ($(p0)!0.52!(p3)$);
  \draw[pair] ($(p1)+(0,0.22)$) -- ($(p1)!0.52!(p2)$);
  \draw[pair] ($(p2)+(-0.22,0)$) -- ($(p2)!0.52!(p3)$);
  \draw[pair] ($(p0)!0.55!(p1)$) -- ++(0,0.62);
  \node[mvtx] at (p0) {0}; \node[mvtx] at (p1) {1};
  \node[mvtx] at (p2) {2}; \node[cvtx] at (p3) {3};
  \node[lab, below] at ($(p0)!0.5!(p1)+(0,-0.02)$) {$01$};
  \node[lab, right] at ($(p1)!0.78!(p2)$) {$12$};
  \node[lab, above] at ($(p2)!0.78!(p3)$) {$23$};
  \node[lab, left] at ($(p0)!0.78!(p3)$) {$03$};
  \node[lab, above left=-1pt] at ($(p0)!0.42!(p2)$) {$02$};
  \node[lab] at (1.62,0.42) {$012$};
  \node[font=\scriptsize] at (1.05,-0.62) {$\beta=(1,1,0)$};
\end{tikzpicture}
\end{minipage}\hfill
%
% ====================== (b) collapse and order ============================
\begin{minipage}[t]{0.68\textwidth}
\textbf{(b)} Collapse and Morse order\par\medskip
{\scriptsize Collapse steps: pair $(01,012)$, pair $(1,12)$; no free face is left,
so $02$ is critical; pair $(2,23)$, pair $(0,03)$; $3$ is critical.}\par\smallskip
\begin{tabular}{@{}r@{\hspace{0.9em}}*{10}{c}@{}}
cell & $\mathbf{3}$ & $0$ & $03$ & $2$ & $23$ & $\mathbf{02}$ & $1$ & $12$ & $01$ & $012$\\
\cmidrule(lr){3-4}\cmidrule(lr){5-6}\cmidrule(lr){8-9}\cmidrule(lr){10-11}
& crit. & \multicolumn{2}{c}{pair} & \multicolumn{2}{c}{pair} & crit.
  & \multicolumn{2}{c}{pair} & \multicolumn{2}{c}{pair}\\[3pt]
key $\ord$ & 0 & 1 & 2 & 3 & 4 & 5 & 6 & 7 & 8 & 9\\
removal time $\rmt$ & 9 & 8 & 7 & 6 & 5 & 4 & 3 & 2 & 1 & 0
\end{tabular}
\end{minipage}

\bigskip
%
% ====================== (c) boundary matrices =============================
\begin{minipage}[t]{0.5\textwidth}
\textbf{(c)} Boundary matrices in the Morse order\par\medskip
\centering
$B_2=\begin{array}{@{}r|c@{}}
    & 012\\ \hline
    03 & 0\\
    23 & 0\\
    02 & 1\\
    12 & 1\\
    01 & \pv{1}\end{array}$\qquad
$B_1=\begin{array}{@{}r|ccccc@{}}
    & 03 & 23 & \cc 02 & 12 & \gz{01}\\ \hline
    3 & 1 & 1 & \cc 0 & 0 & \gz{0}\\
    0 & \pv{1} & 0 & \cc 1 & 0 & \gz{1}\\
    2 & 0 & \pv{1} & \cc 1 & 1 & \gz{0}\\
    1 & 0 & 0 & \cc 0 & \pv{1} & \gz{1}\end{array}$\par\smallskip
\raggedright{\scriptsize Boxed: the apparent pivots, $R_y=\bd y$ with low
$x$ for $(x,y)\in M$. Column $01$ (grey) is cleared, since $01$ is the low of
$R_{012}$. Only the shaded critical column $02$ is reduced
(\cref{cor:critical-work}).}
\end{minipage}\hfill
%
% ====================== (d) reduction of the critical column ===============
\begin{minipage}[t]{0.46\textwidth}
\textbf{(d)} Reducing the critical column $02$\par\medskip
\centering
$\begin{array}{@{}r|c@{}}
    & \bd 02\\ \hline
    3 & 0\\
    0 & 1\\
    2 & \pv{1}\\
    1 & 0\end{array}
\xrightarrow{\;\mu(2)=23\;}
\begin{array}{@{}c@{}}
    {+}\bd 23\\ \hline
    1\\
    \pv{1}\\
    0\\
    0\end{array}
\xrightarrow{\;\mu(0)=03\;}
\begin{array}{@{}c@{}}
    {+}\bd 03\\ \hline
    0\\
    0\\
    0\\
    0\end{array}$\par\smallskip
\raggedright{\scriptsize Each addition replaces the low, a lower cell, by the
other facet of its partner, following the gradient paths $2\to23\to3$ and
$0\to03\to3$ (\cref{cor:morse-reduction}). The column ends at
$\dM 02=3+3=0$: $V_{02}=02+23+03$ is the cycle around the empty triangle, so
$02$ is positive.}
\end{minipage}}%

\caption{The Morse order on a small complex $K$: the square $0123$ with
diagonal $02$ and the triangle $012$ (shaded); the triangle $023$ is not in $K$.
(a)~Arrows point from $x$ to $\mu(x)$ for the gradient pairs found by the
collapse; the critical cells $3$ and $02$ are drawn heavy.
(b)~Reading the removal times from right to left replays the collapse; the keys
are $\ord=9-\rmt$. Each pair is adjacent in the order and its lower cell is the
youngest facet of its partner, so each pair is apparent (\cref{lem:apparent}).
(c)--(d)~Rows and columns are in key order, rows from top to bottom, so a low is
the lowest non-zero entry (boxed).
Every entry is recomputed and checked by \texttt{reference/figure\_examples.py}.}
\label{fig:morse}
\end{figure}

\paragraph{Morse boundary.} By Lemma~\ref{lem:apparent}(b) the following
recursion is well founded and defines a linear map
$\flow\colon\Cc{k-1}\to\Cc{k-1}$: $\flow(\tau)=\tau$ for $\tau\in\Crit$,
$\flow(\tau)=0$ for $\tau\in\Up$, and
$\flow(\tau)=\sum_{\tau'\in\bd\mu(\tau),\,\tau'\ne\tau}\flow(\tau')$ for
$\tau\in\Lo$. Thus $\flow(\tau)$ is the sum, over all gradient paths from $\tau$
to a critical cell, of their endpoints, counted mod~2. The \emph{Morse boundary}
of a critical $k$-cell $c$ is $\dM c:=\flow(\bd c)$, the boundary of the Morse
complex of the matching~\cite{forman2002user}; the Morse boundary matrices
$\Bm{k}^M$ have rows $\Crit\cap\Kc{k-1}$ and columns $\Crit\cap\Kc{k}$. For
$y\in\Up$ with partner $x$, the definition of $\flow(x)$ gives
\begin{equation}\label{eq:flow-kills}
  \flow(\bd y)=\flow(x)+\sum_{f\in\bd y,\,f\ne x}\flow(f)=0 .
\end{equation}

\begin{corollary}[Critical columns perform the Morse reduction]\label{cor:morse-reduction}
Consider the static reduction with clearing under the Morse order and, next to
it, the standard reduction with clearing of the Morse boundary matrices
$\Bm{d}^M,\dots,\Bm{1}^M$, critical cells taken in increasing key order. For
every critical $k$-cell $c$, $k\ge1$:
\begin{enumerate}
\item[(i)] the additions performed on column $c$ are first additions of apparent
  columns $R_y=\bd y$ ($y\in\Up$), then additions of columns of critical cells;
\item[(ii)] if $c$ is not cleared, then when the first phase ends the column
  contains no cell of $\Lo$ and its restriction to $\Crit$ equals $\dM c$;
\item[(iii)] the second phase performs, one for one, the additions that the
  Morse reduction performs on column $c$; $c$ is cleared exactly when it is
  cleared in the Morse reduction; and the final column $R_c$ contains no cell
  of $\Lo$, while $R_c\cap\Crit$ is the reduced Morse column of $c$.
\end{enumerate}
\end{corollary}

Each first-phase addition $R\gets R+\bd y$ with $\low(R)=x$ and $y=\mu(x)$
replaces $x$ by the other facets of $y$; it is one step of the discrete gradient
flow, and $\flow(R)$ does not change by~\eqref{eq:flow-kills}. This is the
precise sense in which the reduced column of a critical cell ``traces its
gradient flow'' (\cref{fig:morse}d). $R_c$ may still contain cells of $\Up$; we observed this in
the randomised checks and claim nothing about them.

\begin{proof}
We use the rule of Definition~\ref{def:collapse} that a critical cell is chosen
of the highest dimension present in $P$. (Lemma~\ref{lem:apparent} did not use
it.)

\emph{Claim (separation).} For every $q$, every critical $q$-cell has a smaller
key than every $q$-cell of $\Lo$. Indeed, when a critical $q$-cell $t$ is removed,
$P$ contains no $(q+1)$-cell. If $\ell\in\Lo$ were a $q$-cell with
$\rmt(\ell)>\rmt(t)$, its partner $\mu(\ell)$, a $(q+1)$-cell removed at time
$\rmt(\ell)-1\ge\rmt(t)$, would be in $P$ at that moment.

We argue by induction over the dimensions (top-down) and, within $\Bm{k}$, over
the columns in increasing key order. The statement about clearing is shown below;
a cleared column receives no addition and has $R_c=0$, matching the zero column
of the Morse reduction. So let $c$ be a column that is not cleared, and let $R$
denote its column during its reduction; initially $R=\bd c$.

\emph{First phase.} Suppose $\ell=\low(R)\in\Lo$ and let $y=\mu(\ell)$. All cells
added to $V_c$ so far are $c$ and upper cells with smaller keys, so $\ell$ is a
facet of a cell $x$ with $\ord(x)\le\ord(c)$; since $y$ is the oldest coface of
$\ell$ (Lemma~\ref{lem:apparent}(c)) and $y\ne c$, $\ord(y)<\ord(c)$. By
Lemma~\ref{lem:apparent}(e), the column of $y$ has been processed, holds
$R_y=\bd y$ and owns $\pivot(\ell)$; the algorithm adds it, which
strictly decreases the low and preserves $\flow(R)$ by~\eqref{eq:flow-kills}.
The phase ends with $R=0$ or with $t=\low(R)\notin\Lo$. In the latter case $R$
is a non-zero cycle (by (I1), $R=\bd V_c$), so $t\notin\Up$ by
Lemma~\ref{lem:apparent}(d), and $t\in\Crit$. By the separation claim every
cell of $R$, having key at most $\ord(t)$, lies outside $\Lo$. As $\flow$ fixes
critical cells and kills upper cells,
$R\cap\Crit=\flow(R)=\flow(\bd c)=\dM c$. This proves (ii) and the first part
of (i).

\emph{Clearing.} The cells cleared in $\Bm{k}$ are the lows of the non-zero
reduced columns of $\Bm{k+1}$: the lower cells (lows of apparent columns) and
the lows of non-zero critical columns. By the induction hypothesis (iii) for
dimension $k+1$, the latter are critical and coincide with the lows of the
non-zero reduced Morse columns, i.e.\ with the cells cleared in the Morse
reduction. So a critical $c$ is cleared in one reduction if and only if in the
other; in the top dimension nothing is cleared in either.

\emph{Second phase.} A pivot on a critical row cannot be owned by an apparent
column (its low is in $\Lo$) or a cleared one (zero), so it is owned by a
critical column $c'$ processed earlier. By the induction hypothesis $R_{c'}$ has
no cell of $\Lo$, it is a cycle, its low is critical and equals
$\low(R_{c'}\cap\Crit)$, and $R_{c'}\cap\Crit$ is the reduced Morse column of
$c'$; hence the pivot tables of the two reductions agree on critical rows.
Adding such columns keeps $R\cap\Lo=\emptyset$ and keeps $R$ a cycle, so while
$R\ne0$ its low is critical and equals $\low(R\cap\Crit)$. The sequence of pivot
look-ups and additions on $R\cap\Crit$, which starts from $\dM c$, is therefore
exactly the Morse reduction of column $c$, and no apparent column is added any
more. This proves (i) and (iii).
\end{proof}

\begin{remark}[Relation to known results]\label{rem:morse-literature}
Lemma~\ref{lem:apparent} is not new in substance. Its part~(c) is the
elementary fact noted in its proof, and the relation between gradient pairs and
apparent pairs is known in both directions: the apparent pairs of any
simplexwise filtration form a discrete gradient, which Ripser uses to skip
columns~\cite{bauer2021ripser}, and Bauer and Roll~\cite{bauer2024wrapping}
relate the gradient pairs of a (generalised) discrete Morse function to the
zero-persistence apparent pairs of a lexicographically refined filtration
order. Lemma~\ref{lem:apparent} is a self-contained version for the order
produced by our collapse procedure, in the form used below: for every queue
discipline, with and without clearing, and with part~(d), on which
Corollary~\ref{cor:morse-reduction} rests. Mischaikow and Nanda~\cite{mischaikow2013morse} showed that
Morse reductions compatible with a filtration preserve persistent homology;
Corollary~\ref{cor:morse-reduction} is a matrix-level form of this for our order:
on critical rows and columns, the static reduction \emph{is} the reduction of
the Morse complex, and by Corollary~\ref{cor:critical-work} the other columns
cost no additions. This is what the \method{full-morse} baseline exploits. In the
language of algebraic Morse theory~\cite{skoldberg2006morse,jollenbeck2009minimal},
the first phase of
Corollary~\ref{cor:morse-reduction} is Gaussian elimination of the invertible
entries $\langle\bd y,x\rangle=1$ of the matched pairs, which is how the Morse
complex arises algebraically; the order makes the reduction algorithm perform
these eliminations by itself, without a separate Morse-reduction pass.
\end{remark}

\begin{remark}[Scope]\label{rem:morse-scope}
Lemma~\ref{lem:apparent} and its corollaries describe the state produced by a
rebuild. Between rebuilds the matching is not maintained, inserted cells get
keys larger than all others, and deletions add chains whose lows may exceed
those of their targets. We make no claim that the DRB stays sparse between
rebuilds; the Morse order is a heuristic for small footprints and short
reduction chains, its effect is measured, and the rebuild policy bounds the
drift of $\nnz(R)+\nnz(V)$.
\end{remark}
